% Part: propositional-logic
% Chapter: syntax-and-semantics
% Section: introduction

\documentclass[../../../include/open-logic-section]{subfiles}

\begin{document}

\olfileid{pl}{syn}{int}

\olsection{مقدمه}

منطق گزاره‌ای به دسته‌ای از !!{formula}s می‌پردازد که از
!!{propositional variable}s با استفاده از رابط‌های گزاره‌ایِ
$\lnot$، $\land$، $\lor$، $\lif$ و $\liff$ ساخته می‌شوند. از دید
شهودی، !!a{propositional variable}~$p$ نمایندهٔ جمله یا گزاره‌ای است که
صادق یا کاذب است. هرگاه «ارزش صدق» هر !!{propositional variable} در
!!a{formula} تعیین شود، ارزش صدق هر یک از !!{formula}s که با استفاده از
رابط‌های گزاره‌ای از این متغیرها ساخته می‌شوند نیز تعیین می‌شود. می‌گوییم
منطق گزاره‌ای \emph{تابع‌ارزشی} است، زیرا معناشناسی آن بر پایهٔ تابع‌هایی
از ارزش‌های صدق تعیین می‌شود. به‌ویژه، در منطق گزاره‌ای هرگونه تعیینِ بیشترِ
صدق و کذب را از بررسی کنار می‌گذاریم؛ برای نمونه، اینکه چیزی ضرورتاً صادق
است نه فقط به‌نحو غیرضروری، اینکه صادق‌بودنِ چیزی معلوم است، یا اینکه
چیزی اکنون صادق است نه اینکه در گذشته صادق بوده یا در آینده صادق خواهد
بود. فقط دو ارزش صدق، یعنی صادق ($\True$) و کاذب ($\False$)، را در نظر
می‌گیریم و ازاین‌رو این امکان را از بحث کنار می‌گذاریم که گزاره‌ای نه صادق
باشد نه کاذب، یا فقط نیمه‌صادق باشد. همچنین تنها بر رابط‌هایی تمرکز می‌کنیم
که ارزش صدقِ !!a{formula} ساخته‌شده از آن‌ها کاملاً از ارزش‌های صدقِ
اجزایش تعیین می‌شود، نه مثلاً از معنای آن. به‌ویژه، اینکه ارزش صدقِ
شرطی‌های زبان انگلیسی به این معنا تابع‌ارزشی است یا نه، محل اختلاف است.
شرطی مادی~$\lif$ تابع‌ارزشی است؛ منطق‌های دیگر به شرطی‌هایی می‌پردازند که
تابع‌ارزشی نیستند.

برای بسط نظریه و فرانظریهٔ منطق گزاره‌ایِ تابع‌ارزشی، نخست باید نحو و
معناشناسی عبارت‌های آن را تعریف کنیم. یک روش ساختن !!{formula}s از
!!{propositional variable}s با استفاده از رابط‌ها را شرح خواهیم داد.
تعریف‌های جایگزین نیز ممکن‌اند. دستگاه‌های دیگر نمادهای متفاوتی برمی‌گزینند،
مجموعه‌های متفاوتی از رابط‌ها را به‌عنوان رابط‌های اولیه انتخاب می‌کنند و
پرانتزها را به شیوه‌ای متفاوت به کار می‌برند؛ یا حتی، مانند آنچه در
نمادگذاری موسوم به لهستانی رخ می‌دهد، اصلاً پرانتز به کار نمی‌برند. بااین‌حال،
وجه مشترک همهٔ رویکردها این است که قواعد ساخت، مجموعهٔ !!{formula}s را
\emph{به‌طور استقرایی} تعریف می‌کنند. اگر این کار درست انجام شود، هر عبارت
اساساً تنها از یک راه مطابق قواعد ساخت به دست می‌آید. اینکه تعریف استقرایی
به عبارت‌هایی \emph{با خوانش یکتا} می‌انجامد، بدین معناست که می‌توانیم با
همان روش، یعنی تعریف استقرایی، به این عبارت‌ها معنا بدهیم.

معنا بخشیدن به عبارت‌ها در قلمرو معناشناسی است. مفهوم محوری در معناشناسیِ
منطق گزاره‌ای، ارضا در !!a{valuation} است. !!^a{valuation}~$\pAssign{v}$
ارزش‌های صدقِ $\True$ و $\False$ را به !!{propositional variable}s نسبت
می‌دهد. هر !!{valuation} یک ارزش صدق به‌صورت $\pValue{v}(!A)$ را برای هر
!!{formula}~$!A$ تعیین می‌کند. !!^a{formula} در
!!a{valuation}~$\pAssign{v}$ ارضا می‌شود اگر و تنها اگر
$\pValue{v}(!A) = \True$؛ این را به‌صورت $\pSat{v}{!A}$ می‌نویسیم. این
رابطه را می‌توان با استقرا بر ساختار~$!A$ نیز تعریف کرد: با استفاده از
توابع ارزشِ رابط‌های منطقی، برای نمونه، ارضای $!A \land !B$ را برحسب ارضا
شدن یا نشدنِ $!A$ و~$!B$ تعریف می‌کنیم.

سپس بر پایهٔ رابطهٔ ارضای $\pSat{v}{!A}$ برای جمله‌ها می‌توانیم مفاهیم
معناییِ بنیادیِ همان‌گویی، استلزام معنایی و ارضاپذیری را تعریف کنیم.
!!^a{formula} یک همان‌گویی است، $\Entails !A$، اگر هر !!{valuation} آن را
ارضا کند؛ یعنی $\pValue{v}(!A) = \True$ برای هر~$\pAssign{v}$.
آن فرمول پیامدِ معناییِ مجموعه‌ای از !!{formula}s است،
$\Gamma \Entails !A$، اگر هر !!{valuation} که همهٔ !!{formula}s در~$\Gamma$
را ارضا می‌کند،~$!A$ را نیز ارضا کند. همچنین مجموعه‌ای از !!{formula}s
ارضاپذیر است اگر یک !!{valuation} همهٔ !!{formula}s آن را هم‌زمان ارضا
کند. چون !!{formula}s به‌طور استقرایی تعریف شده‌اند و ارضا نیز به‌نوبهٔ
خود با استقرا بر ساختار !!{formula}s تعریف می‌شود، می‌توانیم با استقرا
ویژگی‌های معناشناسی خود را ثابت کنیم و میان مفاهیم معناییِ تعریف‌شده
ارتباط برقرار کنیم.


\end{document}
